\section{Methodology}
\label{sec:methodology}

The experimental programme is organized around three complementary workflows.
The chronological phase numbering is retained for reproducibility, but the
scientific interpretation of the experiments is workflow-based:
\begin{description}[leftmargin=!,labelwidth=4.8cm]
  \item[Workflow A.]
        \textbf{Angle-structured QAE integration.}
        A function \(g:[0,1]\to[0,1]\) is encoded as an amplitude under a
        uniform quadrature rule. The target is the numerical integral
        \(I[g]\), approximated by left, midpoint, right, or Simpson quadrature
        rules and then estimated through QAE-type circuits.
  \item[Workflow B.]
        \textbf{Grover--Rudolph probability-law preparation.}
        A positive profile \(F\) defines a finite probability law
        \(p_i=F(x_i)/\sum_jF(x_j)\) on a dyadic grid. The hardware prepares
        \(\sum_i\sqrt{p_i}\ket{i}\), measures the resulting law, and computes
        finite-law quantities of interest
        \(I_p(f)=\sum_i f(x_i)p_i\).
  \item[Workflow C.]
        \textbf{Backend-aware candidate selection.}
        Circuits from Workflows A and B are admitted to QPU execution only
        after local construction, transpilation audit, ideal simulation,
        backend-noise simulation, layout inspection, mitigation selection, and
        explicit rejection tests.
\end{description}

This separation is essential. Workflow A is the direct continuation of the
angle-structured numerical-integration model of
\cite{chinesta2026encoding}. Workflow B is a different experimental object:
it validates Grover--Rudolph preparation of a finite probability law and then
uses that law for direct algebraic integration. Workflow C is not a separate
mathematical estimator; it is the control layer that decides which circuits are
credible enough to spend IBM Quantum credits.

\subsection{Workflow A: angle-structured QAE integration}
\label{subsec:workflow-a}

For Workflow A, the selected quadrature rule \(r\in\{\mathrm{L},\mathrm{M},\mathrm{R}\}\)
fixes the grid points \(x_i^{r}\) introduced in
Section~\ref{subsec:dyadic-grid-notation}. The amplitude circuit \(A_g^{r}\)
prepares
\[
  A_g^{r}\ket{0}^{\otimes n}\ket{0}
  =
  \frac{1}{\sqrt{N}}\sum_{i=0}^{N-1}
  \ket{i}
  \left(
    \sqrt{1-g(x_i^{r})}\ket{0}
    +
    \sqrt{g(x_i^{r})}\ket{1}
  \right),
  \qquad N=2^n.
\]
Measuring the final qubit gives
\[
  P(1;A_g^{r})
  =
  Q_n^{r}[g]
  =
  \frac{1}{N}\sum_{i=0}^{N-1} g(x_i^{r}),
\]
which is a quadrature estimate of the integral of \(g\). Simpson quadrature is
treated as the classical weighted combination of the corresponding left,
midpoint, and right estimates, not as a separate support set \(D_n\). The
current calibration family \(g_0,g_1,g_2\) belongs to the first
angle-structure classes studied in \cite{chinesta2026encoding}. The planned
quadrature extension will separate the discretization error of the rule from
the statistical and hardware error of the QAE estimator.

\subsection{Workflow B: Grover--Rudolph finite-law integration}
\label{subsec:workflow-b}

For Workflow B, the function used to build the quantum state is not directly
the test function to be integrated. A positive profile \(F\) defines the finite
law
\[
  p^{(n)}_j
  =
  \frac{F(x_j)}{\sum_{\ell=0}^{2^n-1}F(x_\ell)},
  \qquad
  x_j=\frac{j+1/2}{2^n},
\]
and the Grover--Rudolph circuit prepares
\[
  \ket{\psi_p}=\sum_{j=0}^{2^n-1}\sqrt{p^{(n)}_j}\ket{j}.
\]
Once the law has been measured, any classical test function \(f\) on the same
grid defines the finite algebraic quantity
\[
  I_{p^{(n)}}(f)=\sum_{j=0}^{2^n-1}f(x_j)p^{(n)}_j.
\]
This is the interpretation used in the distribution, finite-law integration,
MC/QMC baseline, scaling, and Gaussian-profile experiments.

\subsection{Workflow C: backend-aware admission rules}
\label{subsec:workflow-c}

Workflow C controls when hardware execution is allowed. It records logical and
transpiled circuits, backend identity, physical layout, depth, two-qubit-gate
count, two-qubit depth, shots, seeds, expected ideal quantities, ideal
simulation, and backend-noise simulation. A candidate is rejected before QPU
execution when its depth, two-qubit count, noisy total-variation distance, or
largest predicted quantity-of-interest error exceeds the current phase
thresholds. The same workflow also records mitigation choices such as
dynamical decoupling and randomized compiling/twirling.

\subsection{Phase 0: Reproducible audit before spending credits}

Phase 0 generates the calibration circuits parametrically from a single Qiskit
implementation. The circuit generator follows the dyadic-controlled-rotation
logic of Grover--Rudolph state preparation \cite{falco2026grproof}, specialized
here to the small angle-polynomial test family motivated by
\cite{chinesta2026encoding}. For each case \(g_i\) and amplification level \(k\),
the audit stores:
\begin{itemize}[leftmargin=*]
  \item the logical circuit and, when a backend is selected, the transpiled
        circuit;
  \item backend name and status;
  \item logical depth, transpiled depth, two-qubit gate count, two-qubit depth,
        and operation counts;
  \item structural metrics of the encoding;
  \item shot count, seed, and expected ideal probability \(p_k(a)\);
  \item ideal simulation and backend-noise simulation when the corresponding
        simulator is available.
\end{itemize}

The audit rejects candidates whose transpiled depth or two-qubit-gate count
exceeds predefined thresholds. In this first campaign all \(g_0,g_1,g_2\) cases
with \(k=0,1,2\) were below the risk thresholds on \texttt{ibm\_fez}.

\subsection{Larger Grover--Rudolph state preparation}

The next audit extends the workflow from the three-qubit calibration circuits
to genuine Grover--Rudolph loading circuits on larger dyadic grids. Following
the finite probability model in
Section~\ref{subsec:algebraic-probability-model}, each benchmark starts from a
profile \(F(x_i)\) on the midpoint grid \(D_n\), normalizes it to a probability
vector \(p_i\), and encodes it as
\[
  \sum_{j=0}^{2^n-1}\sqrt{p_j}\ket{j}.
\]
Following the probability-tree construction of \cite{falco2026grproof}, for
each binary prefix \(b\) the generator computes the two child masses \(P(b0)\)
and \(P(b1)\), then applies a conditional rotation satisfying
\[
  \sin^2(\theta_b/2)=\frac{P(b1)}{P(b0)+P(b1)}.
\]
The implementation groups all rotations at a fixed tree level into one
uniformly controlled \(R_y\) stage. In Qiskit this is represented with
\texttt{UCRYGate}, matching the ancilla-free Gray-code/Walsh--Hadamard ladder
viewpoint of \cite{falco2026grproof}. This is materially different from
compiling every prefix as an independent multi-controlled rotation: on the
first \(n=4\) audit it reduced the transpiled two-qubit count from hundreds of
gates to 16.

\subsection{Integration scaling from finite probability laws}

After the direct-distribution and cross-backend validation steps, the next
methodological question is not whether a deeper Grover iterate can be run, but
whether algebraic integration remains stable as the dyadic grid is refined. For
a positive profile \(F\) on \([0,1]\), Phase 9 defines the midpoint grid
\[
  x_j=\frac{j+1/2}{2^n},\qquad j=0,\ldots,2^n-1,
\]
and the normalized finite probability law
\[
  p^{(n)}_j
  =
  \frac{F(x_j)}{\sum_{\ell=0}^{2^n-1}F(x_\ell)}.
\]
For each classical test function \(f\), the finite algebraic integral is
\[
  I_{p^{(n)}}(f)
  =
  \sum_{j=0}^{2^n-1} f(x_j)p^{(n)}_j.
\]
This is compared with the continuous reference
\[
  I_F(f)
  =
  \frac{\int_0^1 f(x)F(x)\,dx}{\int_0^1 F(x)\,dx}.
\]
The audit therefore separates three effects:
\begin{align*}
  \epsilon_{\mathrm{disc}}(f)
  &=
  \left|I_{p^{(n)}}(f)-I_F(f)\right|,\\
  \epsilon_{\mathrm{load}}(f)
  &=
  \left|I_{\tilde p_{\mathrm{exact}}}(f)-I_{p^{(n)}}(f)\right|,\\
  \epsilon_{\mathrm{noise}}(f)
  &=
  \left|I_{\tilde p_{\mathrm{noise}}}(f)-I_{p^{(n)}}(f)\right|.
\end{align*}
Here \(\tilde p_{\mathrm{exact}}\) is the probability law extracted from the
exact logical statevector, and \(\tilde p_{\mathrm{noise}}\) is the
backend-noise simulation after transpilation. A hardware campaign is considered
only when the circuit-depth, two-qubit-gate, noisy-TVD, and maximum
quantity-of-interest error thresholds are all satisfied.

\subsection{Sampling baselines under the same finite law}

Phase 10 compares the hardware sampling workflow with classical MC and QMC
sampling after the finite law \(p^{(n)}\) has been fixed. This is a deliberately
controlled comparison: all estimators use the same dyadic law, the same grid,
and the same sample budget \(M\). The QPU estimator is
\[
  I_{\hat p_{\mathrm{QPU}}}(f)
  =
  \sum_{j=0}^{2^n-1}f(x_j)\hat p_{\mathrm{QPU},j},
\]
where \(\hat p_{\mathrm{QPU}}\) is the measured hardware frequency law. The MC
baseline draws \(M\) independent samples \(J_1,\ldots,J_M\sim p^{(n)}\) and
uses
\[
  \hat I_{\mathrm{MC}}(f)
  =
  \frac{1}{M}\sum_{m=1}^M f(x_{J_m}).
\]
The QMC baseline uses randomized shifted van-der-Corput points
\[
  u_m=(v_m+\Delta)\bmod 1,\qquad \Delta\sim U(0,1),
\]
and maps them through the inverse CDF of \(p^{(n)}\):
\[
  J_m=\min\left\{j:\sum_{\ell=0}^j p^{(n)}_\ell\ge u_m\right\},
  \qquad
  \hat I_{\mathrm{QMC}}(f)
  =
  \frac{1}{M}\sum_{m=1}^M f(x_{J_m}).
\]
This QMC baseline is intentionally strong because it assumes the exact finite
CDF is classically available. It should therefore be read as a reference for
sampling quality once the law is known, not as a complete cost model for
constructing that law. The comparison reports distribution-level errors
(TVD, Hellinger distance, classical fidelity, \(L^2\) distance, and maximum cell
deviation) and quantity-level MC/QMC RMSE against the absolute QPU error.

\subsection{Backend selection}

The visible Heron backends were listed before hardware submission. Among the
initial candidates, \texttt{ibm\_fez} had the lowest average two-qubit and
readout errors while maintaining zero or low queue depth at the time of the
campaign. The first hardware phase was therefore run on one backend only, to
avoid confounding amplification-level comparisons with backend calibration
differences.

\subsection{Compilation as an experimental variable}

The \(g_0,k=2\) anomaly observed in Phase 1 shows that compilation cannot be
treated as a passive preprocessing step. The logical circuit and the saved QASM
were both ideal-correct, but the hardware result changed substantially when the
optimization level was changed. This is consistent with the gate-synthesis
viewpoint of \cite{falco2026liegates}: elementary operations, their embeddings,
and their discretized native-basis realizations are part of the experimental
object being benchmarked. For this reason the accepted Phase 1 evidence reports
optimization level, depth, two-qubit-gate count, and QASM paths as first-class
metadata.

\subsection{Hardware execution and stopping rules}

Hardware execution is intentionally staged. The first submission used only six
circuits: \(g_0\) with \(K=\{0,2\}\), \(g_1\) with \(K=\{0,1\}\), and \(g_2\)
with \(K=\{0,1\}\). Deferred circuits were submitted only after the previous
results had been fetched and analyzed.

Each hardware run stores a \texttt{submitted\_campaign.json} file with job
identifiers, compilation metrics, expected probabilities, and QASM paths. Result
fetching stores raw counts and computes the ancilla success probability
\(\hat{p}_k\). A compact analysis table then reports the absolute deviation from
the ideal probability and a binomial \(z\)-score. The classification used in
this report is \texttt{pass} for \(|z|<4\), \texttt{watch} for
\(4\le |z|<8\) or moderate absolute bias, and \texttt{fail} for \(|z|\ge 8\)
or large absolute bias.

MLAE estimates are computed by maximizing the likelihood of the observed
binomial counts across the selected schedule \(K\).
